% appP_graph_f.tex -- generated by writing/tools/blueprint_gen.py 1.0 (level 0) on 2026-10-14 from sources 20021204b9; do not edit
\providecommand{\bpNum}[1]{\hyperref[#1]{\ref*{#1}}}%
\providecommand{\bpNumA}[1]{\hyperref[#1]{\textcolor{orange!65!black}{\ref*{#1}}}}%
\providecommand{\bpPage}[1]{\pageref*{#1}}%
\providecommand{\bpLab}[1]{{\ttfamily #1}}%
\providecommand{\bpBrk}{\allowbreak}%
\providecommand{\bpT}[2]{\facttagbox[#1]{\scriptsize #2}}%
\providecommand{\bpU}[1]{{\color{gray!70!black}#1}}%
\begin{figure}[p]
\centering
\rotatebox{90}{\begin{tikzpicture}[
  bp node/.style={draw,line width=0.45pt,rounded corners=1pt,font=\tiny\ttfamily,align=center,
                  inner xsep=1.5pt,inner ysep=1.1pt,minimum height=4mm},
  bp R/.style={bp node,fill=OliveGreen!15,draw=OliveGreen!60!black},
  bp A/.style={bp node,fill=RoyalBlue!15,draw=RoyalBlue!60!black},
  bp P/.style={bp node,fill=BurntOrange!15,draw=BurntOrange!70!black},
  bp C/.style={bp node,fill=Plum!15,draw=Plum!70!black},
  bp V/.style={bp node,fill=BrickRed!15,draw=BrickRed!70!black,double,double distance=0.6pt},
  bp cond/.style={dash pattern=on 1.6pt off 1pt},
  bp math/.style={bp node,fill=white,draw=black!60},
  bp def/.style={bp node,fill=white,draw=black!45,line width=0.3pt,rounded corners=0pt},
  bp conj/.style={bp node,ellipse,inner xsep=0pt,inner ysep=0pt,fill=Plum!15,draw=Plum!70!black},
  bp asm/.style={bp node,fill=orange!8,draw=orange!65!black,dash pattern=on 2pt off 1.2pt},
  bp stub/.style={bp node,font=\tiny,fill=gray!12,draw=gray!55,text=black!70},
  bp fold/.style={bp node,font=\tiny,fill=white,draw=gray!55,dash pattern=on 1pt off 1pt,text=black!70},
  bp unit/.style={font=\scriptsize,text=black!70,inner sep=1pt},
  bp sep/.style={draw=gray!40,line width=0.4pt,dash pattern=on 1pt off 1.5pt},
  bp key/.style={font=\scriptsize,anchor=west,inner sep=0pt},
  bp use/.style={draw=gray!60,line width=0.5pt,-{Stealth[length=2.6pt,width=2.2pt]}},
  bp stubuse/.style={draw=gray!45,line width=0.45pt,dash pattern=on 1.2pt off 1pt,-{Stealth[length=2.6pt,width=2.2pt]}},
  bp asmuse/.style={draw=orange!65!black,line width=0.5pt,dash pattern=on 2pt off 1.2pt,-{Stealth[length=2.6pt,width=2.2pt]}},
  bp neck/.style={draw=red!60!black,line width=1.1pt,-{Stealth[length=3.6pt,width=3.2pt]}},
  bp long/.style={opacity=0.5}]
  \node[bp unit] at (1.245,4.125) {第~1~章};
  \node[bp unit] at (6.464,4.125) {第~19~章};
  \node[bp unit] at (10.748,4.125) {第~20~章};
  \node[bp unit] at (15.650,4.125) {第~21~章};
  \node[bp unit] at (19.495,4.125) {附录~N};
  \draw[bp sep] (3.163,3.985) -- (3.163,-3.855);
  \draw[bp sep] (4.488,3.985) -- (4.488,-3.855);
  \draw[bp sep] (8.325,3.985) -- (8.325,-3.855);
  \draw[bp sep] (12.792,3.985) -- (12.792,-3.855);
  \draw[bp sep] (18.706,3.985) -- (18.706,-3.855);
  \node[bp asm] (n1) at (0.000,-0.872) {\detokenize{intro-}\\\detokenize{wightman}};
  \node[bp asm] (n2) at (1.245,-0.489) {\detokenize{intro-}\\\detokenize{composite}};
  \node[bp asm] (n3) at (2.490,-0.304) {\detokenize{intro-}\\\detokenize{af-coeff}};
  \node[bp asm] (n4) at (0.000,0.020) {\detokenize{intro-}\\\detokenize{bridge}};
  \node[bp R,bp cond] (n5) at (5.277,0.506) {\detokenize{excl-soft-}\\\detokenize{cancel}};
  \node[bp A] (n6) at (6.522,0.462) {\detokenize{excl-}\\\detokenize{pion-ff}};
  \node[bp A] (n7) at (5.277,-0.385) {\detokenize{excl-}\\\detokenize{transition}};
  \node[bp P] (n8) at (7.652,0.503) {\detokenize{excl-}\\\detokenize{endpoint}};
  \node[bp A] (n9) at (5.277,-1.156) {\detokenize{excl-dvcs}};
  \node[bp A] (n10) at (6.522,-0.429) {\detokenize{excl-}\\\detokenize{meson-L}};
  \node[bp V] (n11) at (7.652,-0.414) {\detokenize{excl-}\\\detokenize{meson-T}};
  \node[bp P] (n12) at (5.277,-1.926) {\detokenize{excl-near-}\\\detokenize{threshold}};
  \node[bp A] (n13) at (9.403,-2.201) {\detokenize{heavy-bdpi}};
  \node[bp A] (n14) at (9.403,-1.431) {\detokenize{heavy-}\\\detokenize{btopi-ff}};
  \node[bp P] (n15) at (10.937,-1.329) {\detokenize{heavy-}\\\detokenize{qcdf-lp}};
  \node[bp V] (n16) at (12.093,-1.246) {\detokenize{heavy-}\\\detokenize{qcdf-}\\\detokenize{endpoint}};
  \node[bp A] (n17) at (9.403,-0.416) {\detokenize{heavy-}\\\detokenize{inclusive-}\\\detokenize{b}};
  \node[bp A] (n18) at (9.403,1.859) {\detokenize{heavy-}\\\detokenize{nrqcd-}\\\detokenize{decay}};
  \node[bp conj] (n19) at (9.403,-3.228) {\detokenize{heavy-}\\\detokenize{nrqcd-}\\\detokenize{production}};
  \node[bp A] (n20) at (9.403,0.721) {\detokenize{heavy-}\\\detokenize{quarkonium-}\\\detokenize{fragmentation}};
  \node[bp A] (n21) at (13.754,-0.067) {\detokenize{power-dis-}\\\detokenize{twist4-x}};
  \node[bp P] (n22) at (13.754,-1.081) {\detokenize{power-}\\\detokenize{twist3-}\\\detokenize{sidis}};
  \node[bp conj] (n23) at (15.551,-1.573) {\detokenize{power-}\\\detokenize{twist3-pp}};
  \node[bp R] (n24) at (13.754,0.948) {\detokenize{power-}\\\detokenize{next-}\\\detokenize{to-soft}};
  \node[bp P] (n25) at (15.551,-0.225) {\detokenize{power-nlp-}\\\detokenize{endpoint}};
  \node[bp R] (n26) at (13.754,3.484) {\detokenize{power-}\\\detokenize{rearrangement}};
  \node[bp P] (n27) at (13.754,2.347) {\detokenize{power-}\\\detokenize{adler-}\\\detokenize{large-}\\\detokenize{beta0}};
  \node[bp P] (n28) at (15.551,2.231) {\detokenize{power-}\\\detokenize{renormalon-}\\\detokenize{examples}};
  \node[bp conj] (n29) at (17.547,2.158) {\detokenize{power-}\\\detokenize{renormalon-}\\\detokenize{ht}};
  \node[bp R] (n30) at (13.754,-2.219) {\detokenize{power-}\\\detokenize{ope-}\\\detokenize{convergence}};
  \node[bp P] (n31) at (13.754,-3.356) {\detokenize{power-}\\\detokenize{thooft-}\\\detokenize{model}};
  \node[bp asm] (n32) at (19.495,-0.197) {\detokenize{appN-}\\\detokenize{collins}};
  \node[bp asm] (n33) at (19.495,-1.089) {\detokenize{appN-ward-}\\\detokenize{product}};
  \node[bp asm] (n34) at (19.495,0.694) {\detokenize{appN-}\\\detokenize{endpoint}};
  \node[bp stub] (n35) at (3.826,1.757) {第 2 章};
  \node[bp stub] (n36) at (3.826,2.407) {第 6 章};
  \node[bp stub] (n37) at (3.826,1.107) {第 7 章};
  \node[bp stub] (n38) at (3.826,-0.843) {第 9 章};
  \node[bp stub] (n39) at (3.826,0.457) {第 10 章};
  \node[bp stub] (n40) at (3.826,-2.143) {第 12 章};
  \node[bp stub] (n41) at (3.826,-0.193) {第 13 章};
  \node[bp stub] (n42) at (3.826,-2.793) {第 14 章};
  \node[bp stub] (n43) at (3.826,-1.493) {第 15 章};
  \begin{scope}[on background layer]
  \draw[bp use,bp long] (n3) -- (n17);
  \draw[bp use,bp long] (n3) -- (n21);
  \draw[bp use,bp long] (n4) -- (n6);
  \draw[bp use,bp long] (n4) -- (n7);
  \draw[bp use,bp long] (n4) -- (n9);
  \draw[bp use,bp long] (n4) -- (n10);
  \draw[bp use,bp long] (n4) -- (n13);
  \draw[bp use,bp long] (n4) -- (n14);
  \draw[bp use,bp long] (n4) -- (n17);
  \draw[bp use,bp long] (n4) -- (n18);
  \draw[bp use,bp long] (n4) -- (n20);
  \draw[bp use,bp long] (n4) -- (n21);
  \draw[bp use,bp long] (n4) -- (n22);
  \draw[bp use,bp long] (n4) -- (n25);
  \draw[bp use] (n1) -- (n2);
  \draw[bp use] (n2) -- (n3);
  \draw[bp use,bp long] (n2) -- (n6);
  \draw[bp use,bp long] (n2) -- (n7);
  \draw[bp use,bp long] (n2) -- (n9);
  \draw[bp use,bp long] (n2) -- (n10);
  \draw[bp use,bp long] (n2) -- (n22);
  \draw[bp use,bp long] (n2) -- (n25);
  \draw[bp stubuse,bp long] (n39) -- (n21);
  \draw[bp stubuse,bp long] (n40) -- (n23);
  \draw[bp stubuse,bp long] (n41) -- (n25);
  \draw[bp stubuse,bp long] (n42) -- (n22);
  \draw[bp stubuse,bp long] (n43) -- (n13);
  \draw[bp stubuse,bp long] (n43) -- (n14);
  \draw[bp stubuse,bp long] (n43) -- (n17);
  \draw[bp use] (n5) -- (n6);
  \draw[bp use] (n5) -- (n10);
  \draw[bp use] (n6) -- (n8);
  \draw[bp use] (n10) -- (n11);
  \draw[bp stubuse,bp long] (n35) -- (n24);
  \draw[bp use] (n14) -- (n15);
  \draw[bp use] (n15) -- (n16);
  \draw[bp use] (n22) -- (n23);
  \draw[bp use] (n24) -- (n25);
  \draw[bp use] (n27) -- (n28);
  \draw[bp use] (n28) -- (n29);
  \draw[bp asmuse,bp long] (n32) -- (n5);
  \draw[bp use,bp long] (n32) -- (n7);
  \draw[bp use,bp long] (n32) -- (n9);
  \draw[bp use,bp long] (n32) -- (n21);
  \draw[bp use,bp long] (n32) -- (n22);
  \draw[bp use] (n32) -- (n25);
  \draw[bp asmuse,bp long] (n33) -- (n5);
  \draw[bp use,bp long] (n33) -- (n7);
  \draw[bp use,bp long] (n33) -- (n9);
  \draw[bp use,bp long] (n33) -- (n21);
  \draw[bp use,bp long] (n33) -- (n22);
  \draw[bp use] (n33) -- (n25);
  \draw[bp asmuse,bp long] (n34) -- (n5);
  \draw[bp stubuse,bp long] (n36) -- (n27);
  \draw[bp stubuse,bp long] (n37) -- (n17);
  \draw[bp stubuse,bp long] (n37) -- (n21);
  \draw[bp stubuse] (n38) -- (n6);
  \draw[bp stubuse] (n38) -- (n7);
  \draw[bp stubuse] (n38) -- (n9);
  \draw[bp stubuse] (n38) -- (n10);
  \draw[bp stubuse,bp long] (n38) -- (n21);
  \draw[bp stubuse,bp long] (n38) -- (n22);
  \draw[bp stubuse,bp long] (n38) -- (n25);
  \draw[bp stubuse,bp long] (n38) -- (n33);
  \draw[bp neck] (n32) -- (n10);
  \draw[bp neck] (n34) -- (n6);
  \draw[bp neck] (n34) -- (n10);
  \end{scope}
  \begin{scope}[shift={(-0.524,-4.355)}]
    \node[bp R,anchor=west,minimum width=4.5mm] (k1) at (0,0) {R};
    \node[bp R,bp cond,anchor=west,minimum width=4.5mm,right=0.8mm of k1] (k2) {R};
    \node[bp A,anchor=west,minimum width=4.5mm,right=0.8mm of k2] (k3) {A};
    \node[bp P,anchor=west,minimum width=4.5mm,right=0.8mm of k3] (k4) {P};
    \node[bp C,anchor=west,minimum width=4.5mm,right=0.8mm of k4] (k5) {C};
    \node[bp V,anchor=west,minimum width=4.5mm,right=0.8mm of k5] (k6) {V};
    \node[bp key,right=1.2mm of k6] (t1) {严格性标签（虚线框为 $R_{\rm cond}$）};
    \node[bp conj,anchor=west,minimum width=7mm,minimum height=3.4mm,right=2.5mm of t1] (k7) {};
    \node[bp key,right=1mm of k7] (t2) {猜想};
    \node[bp asm,anchor=west,minimum width=4.5mm,right=2.5mm of t2] (k8) {};
    \node[bp key,right=1mm of k8] (t3) {编号假设};
    \node[bp math,anchor=west,minimum width=4.5mm,right=2.5mm of t3] (k9) {};
    \node[bp key,right=1mm of k9] (t4) {纯数学};
    \node[bp stub,anchor=west,minimum width=4.5mm,right=2.5mm of t4] (k10) {};
    \node[bp key,right=1mm of k10] (t5) {其他单元};
    \node[bp fold,anchor=west,minimum width=4.5mm,right=2.5mm of t5] (k11) {$+k$};
    \node[bp key,right=1mm of k11] {折叠};
    \draw[bp use] (0,-0.42) -- ++(6mm,0) coordinate (e1);
    \node[bp key] (u1) at ($(e1)+(1.2mm,0)$) {依赖（前提指向陈述）};
    \draw[bp asmuse] ($(u1.east)+(2.5mm,0)$) -- ++(6mm,0) coordinate (e2);
    \node[bp key] (u2) at ($(e2)+(1.2mm,0)$) {编号假设 $\to R_{\rm cond}$};
    \draw[bp neck] ($(u2.east)+(2.5mm,0)$) -- ++(6mm,0) coordinate (e3);
    \node[bp key] (u3) at ($(e3)+(1.2mm,0)$) {瓶颈依赖（主表）};
    \draw[bp use,bp long] ($(u3.east)+(2.5mm,0)$) -- ++(6mm,0) coordinate (e4);
    \node[bp key] (u4) at ($(e4)+(1.2mm,0)$) {跨两列以上};
    \draw[bp stubuse] ($(u4.east)+(2.5mm,0)$) -- ++(6mm,0) coordinate (e5);
    \node[bp key] at ($(e5)+(1.2mm,0)$) {来自其他单元};
  \end{scope}
\end{tikzpicture}}
\caption{第五、六部分（第~19--21~章）主要定理的依赖图（由 \texttt{\string\uses} 自动生成）：总表中的陈述与本部分的全部定理、它们在本部分内的直接前提及所用的编号假设（34 个陈述节点、9 个汇总节点）。框的颜色为严格性标签，虚线框为 $R_{\rm cond}$ 条件定理，橙色虚线框为编号假设，椭圆为猜想，双线框为 [V]，白色框为纯数学；灰色框代表本部分之外的单元（该单元中被直接使用的陈述合并为一个节点），虚线白框“$+k$”为按折叠规则省略的 $k$ 个次要陈述。边由前提指向陈述，是可达关系在所画节点上的传递约化（经由未画出陈述的依赖也画成一条边）；橙色虚线边为编号假设到 $R_{\rm cond}$ 陈述的依赖，粗红线为决定全阶论证地位的瓶颈依赖（取自主表“瓶颈”列）。}\label{fig:appP-dependency-f}
\end{figure}
